\section{Methods}
\label{sec:methods}

The study is a simulation campaign followed by a set of experiments on its results. Every case is
one realization of a front-end: its components drawn within manufacturing tolerances, at most one
injected fault, and three electrodes drawn from a population. Two simulations are run per case.
The first takes the measurements that the instrument could take on itself, with the patient's
electrodes connected. The second measures the specifications of the standard, with the test
networks that the standard prescribes. The first gives the features; the second gives the labels.

\subsection{Circuits}
\label{sec:circuits}

Both circuits are single-lead front-ends (lead~I, left arm minus right arm) with the same signal
chain: series protection and radio-frequency filtering at the inputs, bias resistors, an
instrumentation amplifier, a first-order high-pass filter, a gain stage, a second-order
Sallen-Key low-pass filter and a driven-right-leg (DRL) integrator. Table~\ref{tab:circuits}
summarizes them.

\begin{figure*}[!t]
\centering
\includegraphics[width=\textwidth]{schematic_integrated}
\caption{\cA: integrated instrumentation amplifier with its discrete network, on a single
\SI{3.3}{\volt} supply. $V_\mathrm{cal}$, $V_\mathrm{cmt}$ and $I_\mathrm{lo}$ are the internal
self-test sources; nets with the same name are connected. The electrode model at the lower right
is placed between the body and each of the three leads.}
\label{fig:schematic}
\end{figure*}

\cA\ (Fig.~\ref{fig:schematic}) is the object of the study. It is built around a Texas
Instruments INA333~\cite{ina333} on a single \SI{3.3}{\volt} supply with a buffered mid-supply
reference. The values of the discrete network follow from the standard. The first-stage gain is
4.03, so that a $\pm$\SI{300}{\milli\volt} electrode offset stays within the output swing. The
high-pass corner is \SI{0.041}{\hertz} (\SI{1}{\micro\farad}, \SI{3.9}{\mega\ohm}), as required
by the impulse-response test of the standard, which a \SI{0.5}{\hertz} corner does not pass. The low-pass
filter is a Butterworth response at \SI{185}{\hertz}, which leaves the response at
\SI{150}{\hertz} at 0.83 of mid-band. The bias resistors are \SI{10}{\mega\ohm}, the DRL
integrator has unity gain at \SI{1.6}{\kilo\hertz}, and a \SI{100}{\kilo\ohm} resistor limits the
current into the right-leg electrode. The gain-setting resistor is split in two halves whose
midpoint senses the common-mode voltage.

\cB\ replaces the integrated amplifier by three operational amplifiers and seven resistors on
$\pm$\SI{5}{\volt} supplies, keeps the same filters and uses a first-stage gain of 10.36. It has
the resistor symmetries that Chen \emph{et al.}~\cite{chen2025} identify in benchmark circuits
and stands for the circuits of the diagnosis literature. Its schematic is in the repository.

\begin{table}[!t]
\caption{The Two Circuits}
\label{tab:circuits}
\centering
\begin{tabular}{@{}lcc@{}}
\toprule
 & \cA & \cB \\
\midrule
Instrumentation amplifier & INA333 & three op-amps \\
Supply & \SI{3.3}{\volt} & $\pm$\SI{5}{\volt} \\
Gain (first stage / total) & 4.03 / 278 & 10.36 / 414 \\
Resistors / capacitors with faults & 16 / 8 & 20 / 6 \\
Discrete op-amps & 5 & 6 \\
Converter range & 0 to \SI{3.3}{\volt} & $\pm$\SI{2.5}{\volt} \\
Fault conditions & 293 & 307 \\
Simulated cases & \num{63600} & \num{66400} \\
\bottomrule
\end{tabular}
\end{table}

The circuits are simulated with ngspice~\cite{ngspice}, driven from Python. The INA333 is
represented by a behavioral model built from its data sheet: the three-amplifier structure with
its internal resistors, a single-pole response, white input noise, output clamping, and an error
source for the offset and for a finite common-mode rejection ratio (CMRR) of random sign. The
discrete operational amplifiers use a single-pole model with offset, noise and clamping. In a
bench of the amplifier alone, the behavioral model and the manufacturer's macromodel agree in
gain (4.03) and in CMRR at \SI{50}{\hertz} (\SI{102}{\decibel} against \SI{100}{\decibel}). The
macromodel was not used for the campaign: it has fixed typical parameters, which allows neither
Monte Carlo variation of the block nor faults in it, and in the complete circuit it was about
200 times slower and its operating point did not converge to a valid solution.

Healthy variation is modeled by drawing every resistor uniformly within $\pm$\SI{1}{\percent} of
its nominal value and every capacitor within $\pm$\SI{5}{\percent}; amplifier offsets, open-loop
gains and the CMRR of the INA333 are drawn within their data-sheet ranges.

\subsection{Electrodes and Patient}
\label{sec:electrodes}

Each electrode is a half-cell potential in series with a resistance $R_s$ and with a parallel
pair $R_d \parallel C_d$. Half of the cases use gel electrodes and half dry ones. A gel electrode
takes the network that the standard places in series with each lead, \SI{51}{\kilo\ohm} in
parallel with \SI{47}{\nano\farad}, with $R_s=\SI{300}{\ohm}$. A dry electrode takes one of six
materials (conductive polymer, conductive fabric with and without a membrane, stainless steel,
silver, platinum), with the median $R_s$, $R_d$ and $C_d$ fitted by Joutsen \emph{et
al.}~\cite{joutsen2024} converted from an electrode pair to a single electrode; the median $R_d$
ranges from \SI{61}{\kilo\ohm} for stainless steel to \SI{4.5}{\mega\ohm} for fabric. For each
dry case one of the six subjects of their open data is drawn, and that subject's contact
resistance replaces the median $R_d$, with $C_d$ scaled to keep the time constant. Between
subjects it spans \SI{0.26}{\mega\ohm} to \SI{2.8}{\mega\ohm} for the polymer, \SI{0.66}{\mega\ohm}
to \SI{68}{\mega\ohm} for the fabrics, and \SI{0.6}{\kilo\ohm} to \SI{0.58}{\mega\ohm} for the
solid metals. Finally, every parameter of each of the three electrodes is drawn log-uniformly
between half and twice its value, so that imbalance between electrodes is part of the normal
variation. The polymer and the two fabrics are referred to below as porous dry electrodes.

The patient is a body node coupled to the mains and to earth, with the amplifier common isolated
from earth, as in the model of Winter and Webster~\cite{winter1983}.

Two of these choices are assumptions rather than measurements: the gel electrode has no spread
between subjects, and the factor of two between the electrodes of one case is not taken from
data.

\subsection{Specifications and Functional Labels}
\label{sec:specs}

The front-end is treated as part of a diagnostic electrocardiograph. Eleven quantities are
evaluated for every case following clause 201.12.4 of IEC~60601-2-25~\cite{iec60601225}, with the
limits of Table~\ref{tab:specs}. They are measured with the test networks of the standard at the
inputs, not with the patient's electrodes: the CMRR test applies \SI{10}{\volt} rms at 50 and
\SI{60}{\hertz} with \SI{51}{\kilo\ohm} in parallel with \SI{47}{\nano\farad} in each lead in
turn, without and with a $\pm$\SI{300}{\milli\volt} offset; the input-impedance test places
\SI{620}{\kilo\ohm} in parallel with \SI{4.7}{\nano\farad} in series with a lead. A case is
\emph{compliant} when the eleven quantities are within their limits. Both nominal circuits and
500 healthy Monte Carlo realizations of each are compliant.

\begin{table}[!t]
\caption{Specifications Evaluated for Every Case (IEC~60601-2-25, Clause 201.12.4) and Values of
the Nominal Circuits}
\label{tab:specs}
\centering
\setlength{\tabcolsep}{4pt}
\begin{tabular}{@{}lccc@{}}
\toprule
Specification & Limit & \cA & \cB \\
\midrule
Gain error at \SI{10}{\hertz} & $\le$ \SI{5}{\percent} & 0 & 0 \\
Response, \SIrange{0.67}{40}{\hertz}, deviation & $\le$ \SI{10}{\percent} & \SI{0.3}{\percent} & \SI{0.3}{\percent} \\
Response, \SIrange{40}{150}{\hertz}, minimum & $\ge$ 0.70 & 0.83 & 0.83 \\
Response, \SIrange{40}{500}{\hertz}, maximum & $\le$ 1.10 & 1.00 & 1.00 \\
Baseline shift after impulse$^{\mathrm{a}}$ & $\le$ \SI{100}{\micro\volt} & \SI{75}{\micro\volt} & \SI{75}{\micro\volt} \\
Baseline slope after impulse$^{\mathrm{a}}$ & $\le$ \SI{300}{\micro\volt\per\second} & \SI{18}{\micro\volt\per\second} & \SI{18}{\micro\volt\per\second} \\
Common-mode rejection & $\ge$ \SI{89}{\decibel} & \SI{122}{\decibel} & \SI{111}{\decibel} \\
Noise, peak to valley$^{\mathrm{b}}$ & $\le$ \SI{30}{\micro\volt} & \SI{5.4}{\micro\volt} & \SI{3.4}{\micro\volt} \\
Loss with \SI{620}{\kilo\ohm} $\parallel$ \SI{4.7}{\nano\farad} & $\le$ \SI{20}{\percent} & \SI{3.0}{\percent} & \SI{3.0}{\percent} \\
Gain change with $\pm$\SI{300}{\milli\volt} offset & $\le$ \SI{5}{\percent} & 0 & 0 \\
Input range & $\ge$ \SI{5}{\milli\volt} & \SI{5.8}{\milli\volt} & \SI{6.0}{\milli\volt} \\
\bottomrule
\multicolumn{4}{@{}l}{\footnotesize $^{\mathrm{a}}$\SI{3}{\milli\volt}, \SI{100}{\milli\second}
impulse; referred to the input.}\\
\multicolumn{4}{@{}l}{\footnotesize $^{\mathrm{b}}$Referred to the input, \SIrange{0.05}{150}{\hertz};
taken as 6.6 times the rms value.}
\end{tabular}
\end{table}

Each case carries three levels of labels. The \emph{functional} label is compliance, together
with the list of violated specifications and their continuous values. The \emph{location} label
is the component that was altered. The \emph{origin} label is none, circuit or electrode. Because
the specifications belong to the circuit, an electrode fault leaves the case compliant; it
changes the self-test measurements, and the origin label records why. For comparison, the
\emph{percentage} label of the diagnosis literature is also kept: a case is faulty when a
component was altered, whatever the effect.

\subsection{Fault Model}

One fault at most is injected per case (Table~\ref{tab:faults}). Parametric faults include small
deviations on purpose, since deviations of 5 or \SI{10}{\percent} are where functional and
percentage severity disagree most. Each fault condition is simulated 200 times and the healthy
circuit \num{5000} times, with independent draws of tolerances and electrodes.

\begin{table}[!t]
\caption{Fault Model}
\label{tab:faults}
\centering
\setlength{\tabcolsep}{3pt}
\begin{tabular}{@{}p{0.27\columnwidth}p{0.53\columnwidth}cc@{}}
\toprule
Fault & Model & \multicolumn{2}{c}{Conditions} \\
 & & A & B \\
\midrule
Open, short & \SI{1}{\giga\ohm} in series; \SI{1}{\ohm} in parallel & 48 & 52 \\
Parametric & $\pm$5, $\pm$10, $\pm$20, $\pm$\SI{50}{\percent} of nominal & 192 & 208 \\
Capacitor degradation & $-$\SI{20}{\percent} with \SI{10}{\ohm} in series; $-$\SI{50}{\percent} with \SI{100}{\ohm} & 16 & 12 \\
Op-amp & offset of 20 or \SI{50}{\milli\volt}; open-loop gain divided by 100 or 1000 & 20 & 24 \\
Integrated amplifier & offset of 5 or \SI{20}{\milli\volt}; CMRR of 70 or \SI{50}{\decibel}; gain error of $-$5 or $-$\SI{20}{\percent} & 6 & -- \\
Electrode detached & \SI{1}{\giga\ohm} in series (each of three electrodes) & 3 & 3 \\
Electrode degraded & $R_d$ multiplied and $C_d$ divided by 5 or 20 (each electrode, and both input electrodes) & 8 & 8 \\
\bottomrule
\end{tabular}
\end{table}

\subsection{Self-Test Measurements}
\label{sec:measurements}

The instrument is assumed to have three internal test sources (Fig.~\ref{fig:schematic}): a
floating source $V_\mathrm{cal}$ in series with one lead, which gives the \SI{1}{\milli\volt}
calibration pulse and a differential tone; a source $V_\mathrm{cmt}$ at the reference input of
the DRL amplifier, which the body follows and which therefore gives a common-mode tone; and a
differential current $I_\mathrm{lo}$ into the inputs, as in ac lead-off detection, whose response
is proportional to the contact impedance. They are modeled as ideal sources; how the instrument
generates them is outside the scope of the study. All responses are read at the output of the
chain by the converter of the instrument, with the patient's electrodes connected and no ECG or
mains signal present.

Table~\ref{tab:measurements} lists the four sets of measurements. The main feature set has 26
features. The reading of the instrumentation-amplifier output needs one spare converter channel;
a reading of the DRL output would need another and is only used where stated.

\begin{table}[!t]
\caption{Self-Test Measurements}
\label{tab:measurements}
\centering
\setlength{\tabcolsep}{3pt}
\begin{tabular}{@{}lp{0.58\columnwidth}cc@{}}
\toprule
Set & Measurement & Features & Time \\
\midrule
C1 & Dc level at the output and at the instrumentation-amplifier output & 2 & \SI{0.13}{\second} \\
C2 & Differential and common-mode tones at 0.05, 0.5, 10, 50 and \SI{150}{\hertz}: gain, phase and CMRR estimate & 15 & \SI{94}{\second} \\
C3 & Response to the \SI{1}{\milli\volt}, \SI{200}{\milli\second} calibration pulse: baseline, peak, final value, undershoot, tail, rise time, area & 7 & \SI{1}{\second} \\
C4 & Response to a \SI{1}{\nano\ampere} lead-off current at 10 and \SI{100}{\hertz} & 2 & \SI{2}{\second} \\
\bottomrule
\end{tabular}
\end{table}

The simulation results are noise-free. A measurement model is applied when the features are
built, so that it can be varied without new simulations: Gaussian noise of \SI{1}{\milli\volt}
rms at the converter input, a 12-bit converter sampling at \SI{1}{\kilo\hertz}, clipping to its
range, dc readings averaged over 64 samples and tones estimated from \num{1000} samples or two
periods, whichever is longer. Referred to the input of \cA, one code is \SI{2.9}{\micro\volt}. The
two \SI{0.05}{\hertz} tones account for 80 of the \SI{97}{\second} that the complete self-test
takes.

\subsection{Testability Analysis}
\label{sec:testability}

What the measurements can distinguish is analyzed before any model is trained, in two ways.

The first uses only the nominal circuit. The sensitivity of feature $f_i$ to component $p_j$ is
expressed in units of the spread of that feature over healthy circuits,
\begin{equation}
z_{ij} = \frac{0.01\,p_j}{\sigma_i}\,\frac{\partial f_i}{\partial p_j},
\label{eq:sensitivity}
\end{equation}
so that $z_{ij}=1$ means that a \SI{1}{\percent} deviation of the component moves the feature by
one standard deviation of the healthy population. Two components whose columns of
$Z=[z_{ij}]$ are parallel act on the measurements in the same way, and a deviation of one can be
mimicked by a deviation of the other. Components whose columns have an absolute cosine of at
least 0.99 are joined in an \emph{a priori ambiguity group}; the remaining components, the
amplifiers and the electrodes are groups of one. The number of singular values $s$ of $Z$ with
$10\,s\ge 3$ is reported as the testability rank: the number of independent directions along
which a \SI{10}{\percent} deviation moves the features by at least three healthy standard
deviations.

The second uses the simulated faults. Each fault condition is summarized by the median and the
interquartile range of every feature, which, unlike the mean and the standard deviation, are not
distorted by the heavy tails of faults that saturate the output. Two conditions are called
indistinguishable when no single feature separates their medians by three pooled spreads. This
criterion is univariate and therefore conservative; it gives a map of which non-compliant
conditions look healthy and which components can be mistaken for which, against which the trained
models are later compared.

\subsection{Decision, Location and Origin Models}
\label{sec:models}

\subsubsection{Compliance decision}
Four methods decide pass or fail from the same features. The \emph{limit test} is the classical
reference: a case fails if any feature falls outside the interval that contains the healthy
circuits, with the intervals widened so that jointly they keep \SI{99}{\percent} of them. The
\emph{specification regression} follows alternate test: one gradient-boosting regressor per
specification, and a case fails if any predicted value is beyond its limit. The
\emph{classifier} is a gradient-boosting classifier of compliance at its default threshold. The
\emph{classifier with an escape target} lowers the decision threshold to the value that, on the
validation data, lets through \SI{1}{\percent} of the non-compliant cases; the regression has an
equivalent guard band. The quality of the decision is given by the escape rate, the fraction of
non-compliant circuits that are accepted, and by the false-reject rate, the fraction of compliant
circuits that are rejected.

\subsubsection{Functional against percentage severity}
The same classifier, with classes weighted inversely to their frequency, is trained once on the
functional label and once on the percentage label. Both are scored against compliance, which is
what matters to the user.

\subsubsection{Location}
On the non-compliant cases, five classifiers ($k$-nearest neighbors, random forest, support
vector machine, gradient boosting and a multilayer perceptron) are trained to name the a priori
ambiguity group and, separately, the component. A one-dimensional convolutional network on the
\num{1000} samples of the pulse response is included for comparison.

\subsubsection{Origin}
A three-class problem is defined by what has to be done: \emph{none} (a healthy circuit, or a
deviation that leaves it compliant), \emph{circuit} (out of specification) and \emph{electrode}
(an electrode detached or degraded). An open protection resistor in series with a lead is
electrically the same as a detached electrode and calls for the same action, so it belongs to the
electrode class.

The models are those of scikit-learn~\cite{pedregosa2011} with small hyperparameter grids; the
purpose is to establish what the measurements allow, not to optimize a model.

\subsection{Evaluation Protocol}

\subsubsection{Terminology}
In the vocabulary of metrology~\cite{vim,gum}, the eleven specifications are the measurands, and
the self-test estimates them indirectly, or estimates directly whether they are within their
limits. The pass/fail decision is therefore a conformity assessment in the sense of JCGM~106
\cite{jcgm106}. The escape rate is the probability that a non-conforming circuit is accepted,
and the false-reject rate the probability that a conforming one is rejected; they correspond to
the consumer's and the producer's risk, conditioned on the true state of the circuit and
evaluated over the simulated population, whose proportion of non-conforming cases is set by the
fault model and is not that of instruments in service. The escape target acts as a guard band on
the acceptance limit. The quality of the indirect estimate of each specification is reported as
the root-mean-square error of its prediction relative to the limit. Accuracy, balanced accuracy,
recall and F1 are used only as figures of merit of classifiers and are not statements about
measurement accuracy.

\subsubsection{Splits and intervals}
Every experiment is repeated on three random splits into training, validation and test parts of
52.5, 17.5 and \SI{30}{\percent}, stratified so that every fault condition appears in the three
parts. Hyperparameters, thresholds and guard bands are chosen on the validation part only.
Results are means over the three repetitions; intervals are \SI{95}{\percent} confidence
intervals of the mean (Student's $t$). Unless stated otherwise, results use all electrode types
and the main feature set.

Robustness is examined in four ways: holding all parametric faults of some magnitudes out of
training; sweeping the converter noise from 0.1 to \SI{10}{\milli\volt} rms and its resolution
from 8 to 16 bits, with models trained at the swept setting or at the default one; scoring the
models on two additional datasets generated with other seeds, one with truncated Gaussian
tolerances and one with \SI{2}{\percent} resistors and \SI{10}{\percent} capacitors; and raising
the amplitude and duration of the common-mode tone. The smallest useful set of measurements is
found by greedy forward selection over the self-test actions (one dc reading, one tone or the
pulse), adding at each step the action that most improves the validation score.
